\documentclass[10pt]{article}

\usepackage[margin=0.8in]{geometry}
\usepackage{amsmath,amssymb}
\usepackage{booktabs}
\usepackage[hidelinks]{hyperref}

\setlength{\parindent}{0pt}
\setlength{\parskip}{0.4em}

\title{A Transparent Quarterly Model of U.S. Fiscal Dynamics\\
\large Accounting, Treasury Refinancing, and Inflation}
\author{U.S. Fiscal Dynamics Simulator, version 0.1}
\date{August 2026}

\begin{document}
\maketitle

\begin{abstract}
This note describes a deterministic model of U.S. federal debt dynamics.  The model is
designed to answer accounting questions before attempting to model political or
macroeconomic behavior.  Its central feature is a security-level representation of
marketable Treasury debt: fixed-rate securities retain their rates until maturity, bills
roll over quickly, floating-rate notes reset, and Treasury Inflation-Protected Securities
(TIPS) have inflation-indexed principal.  Inflation, real growth, primary deficits, and
new-issuance rates are independent assumptions.  The resulting model is intentionally
simple, but each result can be traced to an equation, an observed Treasury security, a CBO
projection, or a stated approximation.
\end{abstract}

\section{Introduction}

A familiar fiscal projection shows federal debt relative to gross domestic product (GDP)
rising for many years.  Such a projection does not by itself explain how the debt stock
would respond to a sudden change in inflation or interest rates.  Two details are
especially important.  First, unexpected inflation can raise nominal GDP immediately,
lowering the debt-to-GDP ratio even if no nominal debt disappears.  Second, most Treasury
debt does not instantly reprice: higher market rates enter federal interest costs gradually
as securities mature, are refinanced, or reset.

Version 0.1 of the U.S. Fiscal Dynamics Simulator isolates those mechanisms.  It is not a
forecast, a general-equilibrium model, or a model of fiscal crisis.  Congress, the Federal
Reserve, and investors do not react endogenously.  Instead, the user specifies paths for
inflation, real growth, the primary deficit, other financing, and issuance rates.  The
model then calculates the implied GDP, interest costs, borrowing, refinancing, and debt
stock.

The benchmark is the Congressional Budget Office's (CBO's) February 2026 baseline
projection for 2026--2036 \cite{cbo2026outlook}.  The starting Treasury portfolio is the
September 30, 2025 Monthly Statement of the Public Debt (MSPD) \cite{mspd}.

\section{Time, notation, and GDP}

The simulation step is one calendar quarter.  A federal fiscal year contains the December
quarter of the preceding calendar year and the March, June, and September quarters of the
named fiscal year.  CBO budget flows are fiscal-year quantities and are allocated equally
across those four quarters.  CBO quarterly economic data retain their calendar-quarter
timing.

Let $Y_t$ denote nominal GDP at a seasonally adjusted annual rate, $g_t$ annualized real GDP
growth, and $\pi_t$ annualized GDP-price inflation.  Annual rates are converted to effective
quarterly rates:
\begin{equation}
 g^q_t=(1+g_t)^{1/4}-1,
 \qquad
 \pi^q_t=(1+\pi_t)^{1/4}-1.
\end{equation}
Nominal GDP evolves according to
\begin{equation}
 Y_t=Y_{t-1}(1+g^q_t)(1+\pi^q_t).
 \label{eq:gdp}
\end{equation}
This formulation makes inflation and real growth separate, visible assumptions.  The model
does not infer an effect of inflation on real output.

\section{Fiscal and debt accounting}

Let $P_t$ be the primary deficit, with a positive value denoting a deficit and a negative
value denoting a surplus.  Let $I_t$ be modeled debt interest and $F_t$ the explicitly
reported ``other financing adjustment.''  The modeled total deficit is
\begin{equation}
 G_t=P_t+I_t,
 \end{equation}
and debt held by the public, $D_t$, satisfies
\begin{equation}
 D_t=D_{t-1}+G_t+F_t.
 \label{eq:debtidentity}
\end{equation}
The term $F_t$ matters because changes in debt held by the public need not equal the unified
budget deficit.  Treasury cash-balance changes and federal credit-program cash flows are
examples of other means of financing \cite{cbo2026outlook}.  No hidden residual is used to
force equation~\eqref{eq:debtidentity} to match a target series.

The treatment of TIPS makes it useful to distinguish interest expense from immediate cash
borrowing.  Write
\begin{equation}
 I_t=C_t+I^O_t+X_t+L_t,
\end{equation}
where $C_t$ is marketable interest excluding TIPS principal indexation, $I^O_t$ is the
approximated interest cost of other public debt, $X_t$ is TIPS inflation compensation, and
$L_t$ is any TIPS maturity-floor cost.  TIPS indexation raises principal directly.  Thus the
signed cash financing requirement is
\begin{equation}
 B_t=P_t+C_t+I^O_t+F_t,
\end{equation}
while $X_t+L_t$ enters the debt stock through the TIPS liability itself.  Consequently,
$B_t+X_t+L_t=P_t+I_t+F_t$, preserving equation~\eqref{eq:debtidentity} without financing
TIPS compensation twice.  If $B_t<0$, the model reports debt retirement rather than
negative ``new borrowing.''

Quarterly ratios use the annual-rate GDP denominator:
\begin{equation}
 \frac{P_t}{Y_t}\bigg|_{\mathrm{annualized}}=\frac{4P_t}{Y_t},
 \qquad
 \frac{I_t}{Y_t}\bigg|_{\mathrm{annualized}}=\frac{4I_t}{Y_t},
 \qquad
 \frac{D_t}{Y_t}=\frac{D_t}{Y_t}.
\end{equation}

\section{Treasury cohorts and refinancing}

Each marketable Treasury cohort $i$ has a type $k_i$, issue quarter, maturity quarter,
principal $Q_{i,t}$, fixed coupon, and effective financing rate $r_{i,t}$.  The starting
ledger contains 458 CUSIPs.  At the end of quarter $t$, let $\mathcal{M}_t$ be the set of
cohorts that mature.  Maturing principal is
\begin{equation}
 R_t=\sum_{i\in\mathcal{M}_t}Q_{i,t}.
\end{equation}
The model removes $R_t$ and issues the same amount as replacement debt.  This operation
changes the rate and maturity composition but not the total debt stock or the deficit.
Gross Treasury issuance is therefore
\begin{equation}
 \text{gross issuance}_t=R_t+\max(B_t,0),
\end{equation}
whereas genuinely new borrowing is only $\max(B_t,0)$.

\subsection*{Bills and fixed-rate notes and bonds}

Bills are carried at principal with a transparent effective-yield approximation.  Their
quarterly financing cost is $Q_{i,t}r_{i,t}/4$.  Their short representative maturity makes
bill costs respond rapidly to new issuance rates.

For notes and bonds, $r_{i,t}$ is fixed until maturity.  A rate shock at time $t$ therefore
does not change the current cost of a long-dated inherited security.  It affects that
principal only after the security enters $\mathcal{M}_t$ and is replaced.  This is the main
source of gradual repricing in the model.  Starting effective rates use issue-amount-weighted
auction yields when available, otherwise coupons; premium and discount amortization are not
modeled separately.

\subsection*{TIPS}

For a TIPS cohort, let $\pi^{q,\mathrm{ref}}_t$ be the quarterly reference-CPI change.  The
indexed principal and inflation compensation are
\begin{align}
 Q_{i,t}^{+}&=Q_{i,t}^{-}(1+\pi^{q,\mathrm{ref}}_t),\\
 X_{i,t}&=Q_{i,t}^{+}-Q_{i,t}^{-}.
\end{align}
The real coupon is applied to indexed principal, so quarterly coupon cost is
$Q_{i,t}^{+}r_i/4$.  At maturity, redemption is the greater of indexed principal and
original par.  These conventions follow Treasury's TIPS terms and federal budget treatment,
under which principal adjustments and coupon accruals are recorded as interest
\cite{tips,cbonetinterest}.  The official index is daily and uses a lagged reference CPI;
the quarterly model approximates it with a one-quarter lag.

\subsection*{Floating-rate notes}

Treasury floating-rate notes (FRNs) use the most recent 13-week bill auction rate plus a
fixed auction spread and reset weekly \cite{frn}.  At quarterly resolution, the simulator
uses
\begin{equation}
 r^{\mathrm{FRN}}_{i,t}=r^{\mathrm{short}}_t+s_i
\end{equation}
and resets all FRNs once per quarter.  FRN interest therefore responds to a rate shock
without waiting for final maturity.

\subsection*{New issuance}

New deficit borrowing is allocated using the starting public marketable composition:
\begin{center}
\begin{tabular}{lrr}
\toprule
Instrument & Issuance share & Representative tenor \\
\midrule
Bills & 21.54\% & 1 quarter \\
Notes & 51.82\% & 7 years \\
Bonds & 17.29\% & 20 years \\
TIPS & 7.03\% & 10 years \\
FRNs & 2.33\% & 2 years \\
\bottomrule
\end{tabular}
\end{center}
This fixed mix is an explicit approximation, not a prediction of Treasury debt-management
policy.  Short, intermediate, long, TIPS-real, and FRN-spread assumptions are separately
visible.

\section{Debt coverage and interest definitions}

The headline stock is federal debt held by the public, not gross federal debt.  It is
decomposed as
\begin{equation}
 D_t=M_t+O_t,
\end{equation}
where $M_t$ is modeled marketable debt and $O_t$ is other debt held by the public.  MSPD
security detail reports total amounts outstanding rather than public ownership by CUSIP.
The model therefore scales CUSIP principal within each instrument class to the MSPD
public-held class total.  On September 30, 2025, the reconciliation is
\begin{align*}
 D_0 &= \$30{,}277.766\text{ billion},\\
 M_0 &= \$29{,}695.002\text{ billion},\\
 O_0 &= \phantom{\$30{,}}\$582.764\text{ billion}.
\end{align*}

Modeled debt interest is not relabeled as CBO net interest.  CBO net interest consists
primarily of interest on debt held by the public, offset by interest income and adjusted for
other federal budget accounts \cite{cbonetinterest}.  The simulator retains CBO's series as
a comparison and reports the difference.  CBO itself uses security-level principal,
maturity, rate, financing needs, and issuance-mix assumptions when projecting net interest
\cite{cbomethods}.

\section{Scenarios and interpreting inflation}

The baseline uses CBO paths for primary deficits, other financing, real GDP, inflation,
3-month bill rates, and 10-year note rates.  Experimental scenarios may change inflation,
real growth, the primary deficit, and issuance rates independently.  This permits
deliberately artificial comparisons---for example, high inflation with unchanged rates---to
isolate accounting channels without pretending that such a combination is a forecast.

For a baseline $b$ and scenario $s$, the debt-ratio difference is
\begin{equation}
 \Delta d_t=\frac{D^s_t}{Y^s_t}-\frac{D^b_t}{Y^b_t}.
\end{equation}
The simulator uses an exact two-factor Shapley decomposition.  The nominal-debt contribution
is
\begin{equation}
 \Delta d_t^D=\frac{1}{2}\left[
 \left(\frac{D^s_t}{Y^b_t}-\frac{D^b_t}{Y^b_t}\right)
 +\left(\frac{D^s_t}{Y^s_t}-\frac{D^b_t}{Y^s_t}\right)\right],
\end{equation}
and the nominal-GDP contribution is
\begin{equation}
 \Delta d_t^Y=\frac{1}{2}\left[
 \left(\frac{D^b_t}{Y^s_t}-\frac{D^b_t}{Y^b_t}\right)
 +\left(\frac{D^s_t}{Y^s_t}-\frac{D^s_t}{Y^b_t}\right)\right].
\end{equation}
By construction, $\Delta d_t^D+\Delta d_t^Y=\Delta d_t$.  This is accounting attribution,
not a causal estimate.  It makes the central interpretation explicit: inflation may lower
debt/GDP through $Y_t$ while nominal debt $D_t$ continues to rise because of primary
deficits, TIPS compensation, and refinancing costs.

\section{Limitations}

The model has no endogenous response of taxes, program spending, monetary policy, investor
demand, exchange rates, or real economic activity.  It does not model default, crisis
probabilities, or a debt threshold.  Federal Reserve holdings are included in debt held by
the public, but Treasury and Federal Reserve balance sheets are not consolidated; reserve
interest, remittances, and the Fed's portfolio are outside the model.  Other public debt has
no security-level maturity structure, fiscal-year flows are spread evenly across quarters,
and new issuance uses representative rather than auction-specific maturities.  Results are
therefore conditional accounting experiments and should always be reported together with
their imposed inflation, rate, growth, and primary-deficit paths.

\begin{thebibliography}{9}

\bibitem{cbo2026outlook}
Congressional Budget Office.
\emph{The Budget and Economic Outlook: 2026 to 2036}.
February 2026.
\url{https://www.cbo.gov/publication/62105}.

\bibitem{cbonetinterest}
Congressional Budget Office.
\emph{Federal Net Interest Costs: A Primer}.
December 2020.
\url{https://www.cbo.gov/publication/56910}.

\bibitem{cbomethods}
Congressional Budget Office.
\emph{CBO Explains How It Develops the Budget Baseline}.
April 2023.
\url{https://www.cbo.gov/publication/59085}.

\bibitem{mspd}
U.S. Department of the Treasury, Bureau of the Fiscal Service.
\emph{Monthly Statement of the Public Debt}.
September 30, 2025 observation.
\url{https://fiscaldata.treasury.gov/datasets/monthly-statement-public-debt/}.

\bibitem{tips}
U.S. Department of the Treasury, TreasuryDirect.
\emph{Treasury Inflation-Protected Securities (TIPS)}.
\url{https://www.treasurydirect.gov/marketable-securities/tips/}.

\bibitem{frn}
U.S. Department of the Treasury, TreasuryDirect.
\emph{Floating Rate Notes}.
\url{https://www.treasurydirect.gov/marketable-securities/floating-rate-notes/}.

\end{thebibliography}

\end{document}
